\section{Enumeration, pattern lifting, and constrained systems}

The number of de Bruijn cycles of order \(n\) over an alphabet of size \(q\) is
\[
N(q,n)=\frac{(q!)^{q^{n-1}}}{q^n}.
\]
For \(q=2\), this reduces to
\[
N(2,n)=2^{2^{n-1}-n}.
\]
This count follows from the BEST theorem applied to \(B(q,n-1)\).

In the homological framework, an Eulerian circuit is a full-support integral 1-cycle:
\[
\partial c=0.
\]
Hence \(N(q,n)\) counts global cyclic sections of the word sheaf satisfying the full-coverage condition.

The quotient-sheaf viewpoint refines this count. Let
\[
\overline B(q,n)=B(q,n)/S_q
\]
be the equality-pattern quotient. Then
\[
N(q,n)=\sum_S |\operatorname{Lift}(S)|,
\]
where \(S\) ranges over pattern skeletons and \(\operatorname{Lift}(S)\) is the set of compatible symbol assignments in the fiber. Nonzero obstruction monodromy gives
\[
|\operatorname{Lift}(S)|=0.
\]

For constrained systems \(\mathcal C\), such as DNA/RNA assembly constraints, patient-specific treatment constraints, or atmospheric disturbance constraints, the analogous count is
\[
N_{\mathcal C}(q,n)=
\sum_{\text{valid }S}
|\operatorname{Lift}_{\mathcal C}(S)|.
\]
If the constrained graph remains Eulerian, BEST can still be applied. Otherwise,
\[
\partial c\neq0
\]
records the near-\(G\) boundary defects.
